% ============================================================================
%  Chapitre 2, partie I : Mécanique du solide
%  (manuscrit P1, p. 13 à 19)
% ============================================================================

% ----------------------------------------------------------------------------
\section{Mécanique du solide}
% ----------------------------------------------------------------------------

On considère un solide $\mathcal{B}$ de masse $m$, de centre de masse
$C$, et un point courant $P\in\mathcal{B}$ de masse élémentaire
$\dd m$.

% ----------------------------------------------------------------------------
\subsection{Centre de masse et cinématique}

\subsubsection{Centre de masse}

\begin{definition}[Centre de masse]
    Le \emph{centre de masse} $C$ est défini par
    \[
        \overrightarrow{OC}
        = \frac{1}{m}\int_{\mathcal{B}} \overrightarrow{OP}\,\dd m,
        \qquad m = \int_{\mathcal{B}} \dd m
    \]
\end{definition}

En écrivant $\overrightarrow{OP} = \overrightarrow{OC}
+ \overrightarrow{CP}$,
\[
    \overrightarrow{OC}
    = \frac{1}{m}\int_{\mathcal{B}} \overrightarrow{OC}\,\dd m
    + \frac{1}{m}\int_{\mathcal{B}} \overrightarrow{CP}\,\dd m
    = \overrightarrow{OC}
    + \frac{1}{m}\int_{\mathcal{B}} \overrightarrow{CP}\,\dd m
    \longleftrightarrow
    \int_{\mathcal{B}} \overrightarrow{CP}\,\dd m = \vect{0}
\]

\subsubsection{Champ des vitesses d'un solide}

Soient deux points $A$ et $B$ fixes dans un repère $\rep{j}$ en
mouvement par rapport à un repère $\rep{i}$ d'origine $O_i$. La vitesse
de $B$ par rapport à $\rep{i}$ est par définition
\[
    \vect{v}_{B/i}
    = \left(\ddt\overrightarrow{O_iB}\right)_{\rep{i}}
    = \left(\ddt\overrightarrow{O_iA}\right)_{\rep{i}}
    + \left(\ddt\overrightarrow{AB}\right)_{\rep{i}}
    = \vect{v}_{A/i}
    + \left(\ddt\overrightarrow{AB}\right)_{\rep{i}}
\]

\begin{theoreme}[Dérivation en base mobile]
    Soient deux repères $\rep{i}$ et $\rep{j}$, $\rep{j}$ étant en
    mouvement par rapport à $\rep{i}$ avec la vitesse angulaire
    $\vect{\Omega}_{j/i}$. Pour deux points $A$ et $B$,
    \[
        \left(\ddt\overrightarrow{AB}\right)_{\rep{i}}
        = \left(\ddt\overrightarrow{AB}\right)_{\rep{j}}
        + \vect{\Omega}_{j/i}\wedge\overrightarrow{AB}
    \]
\end{theoreme}

Si $A$ et $B$ sont fixes dans $\rep{j}$, $\overrightarrow{AB}$ y est
constant et
\[
    \vect{v}_{B/i} = \vect{v}_{A/i}
    + \vect{\Omega}_{j/i}\wedge\overrightarrow{AB}
\]
Appliquée au solide, avec $\vect{\Omega}$ sa vitesse angulaire, cette
relation donne le \emph{champ des vitesses}
\[
    \vect{v}_P = \vect{v}_C + \vect{\Omega}\wedge\overrightarrow{CP}
\]

\subsubsection{Champ des accélérations}

En dérivant le champ des vitesses dans $\rep{i}$,
\[
    \vect{a}_P = \dot{\vect{v}}_P
    = \dot{\vect{v}}_C
    + \dot{\vect{\Omega}}\wedge\overrightarrow{CP}
    + \vect{\Omega}\wedge
    \left(\ddt\overrightarrow{CP}\right)_{\rep{i}}
\]
et, puisque $\bigl(\ddt\overrightarrow{CP}\bigr)_{\rep{i}}
= \vect{\Omega}\wedge\overrightarrow{CP}$,

\begin{proposition}[Champ des accélérations]
    \[
        \vect{a}_P = \vect{a}_C
        + \dot{\vect{\Omega}}\wedge\overrightarrow{CP}
        + \vect{\Omega}\wedge\bigl(\vect{\Omega}\wedge
        \overrightarrow{CP}\bigr)
    \]
\end{proposition}

% ----------------------------------------------------------------------------
\subsection{Quantité de mouvement et loi de Newton}

\begin{definition}[Quantité de mouvement du solide]
    \[
        \vect{P} = \int_{\mathcal{B}} \vect{v}_P\,\dd m
    \]
\end{definition}

Avec le champ des vitesses et $\int_{\mathcal{B}}\overrightarrow{CP}\,
\dd m = \vect{0}$, on obtient $\vect{P} = m\,\vect{v}_C$

\begin{theoreme}[Loi de Newton]
    Dans un repère galiléen, si $\vect{F}_{\mathrm{ext}}$ est la
    résultante des actions extérieures,
    \[
        \ddt\vect{P} = \vect{F}_{\mathrm{ext}},
        \qquad\text{donc}\qquad
        m\,\vect{a}_C = \vect{F}_{\mathrm{ext}}
    \]
\end{theoreme}

% ----------------------------------------------------------------------------
\subsection{Tenseur d'inertie et moment cinétique}

\begin{definition}[Moment cinétique]
    Le \emph{moment cinétique} au centre de masse est
    \[
        \vect{L}_C = \int_{\mathcal{B}}
        \overrightarrow{CP}\wedge\vect{v}_P\,\dd m
    \]
\end{definition}

En posant $\vect{r} = \overrightarrow{CP}$, le champ des vitesses donne
\[
    \vect{L}_C = \int_{\mathcal{B}}
    \vect{r}\wedge(\vect{\Omega}\wedge\vect{r})\,\dd m,
    \qquad\text{car}\qquad
    \int_{\mathcal{B}} \vect{r}\wedge\vect{v}_C\,\dd m
    = \left(\int_{\mathcal{B}}\vect{r}\,\dd m\right)\wedge\vect{v}_C
    = \vect{0}
\]

\begin{proposition}[Double produit vectoriel]
    \[
        \vect{r}\wedge(\vect{\Omega}\wedge\vect{r})
        = \|\vect{r}\|^2\vect{\Omega} - \vect{r}\,(\vect{r}^T\vect{\Omega})
        = \bigl(\|\vect{r}\|^2\Id - \vect{r}\vect{r}^T\bigr)\vect{\Omega}
    \]
\end{proposition}

On obtient alors
\[
    \vect{L}_C = \left(\int_{\mathcal{B}}
    \bigl(\|\vect{r}\|^2\Id - \vect{r}\vect{r}^T\bigr)\dd m\right)
    \vect{\Omega}
\]

\begin{definition}[Tenseur d'inertie autour de $C$]
    \[
        \mat{I}_C = \int_{\mathcal{B}}
        \bigl(\|\vect{r}\|^2\Id - \vect{r}\vect{r}^T\bigr)\dd m,
        \qquad
        \vect{L}_C = \mat{I}_C\,\vect{\Omega}
    \]
\end{definition}

\begin{theoreme}[Propriétés du tenseur d'inertie]
    \leavevmode
    \begin{itemize}
        \item $\mat{I}_C$ est symétrique : $\mat{I}_C^T = \mat{I}_C$ ;
        \item $\mat{I}_C$ est positive : pour tout $\vect{u}\in\R^3$,  c'est-à-dire
            \[
                \vect{u}^T\mat{I}_C\,\vect{u}
                = \int_{\mathcal{B}}\|\vect{r}\wedge\vect{u}\|^2\,\dd m
                \geq 0
            \]
    \end{itemize}
    Par le théorème spectral, $\mat{I}_C$ est donc diagonalisable en
    base orthonormée : dans la base des \emph{axes principaux
    d'inertie},
    \[
        \mat{I}_C = \mathrm{diag}(I_1,I_2,I_3),
        \qquad I_1, I_2, I_3 \geq 0
    \]
\end{theoreme}

% ----------------------------------------------------------------------------
\subsection{Théorème du moment cinétique}

On note $\vect{M}_C^{\mathrm{ext}}$ le moment résultant des actions
extérieures autour de $C$.

\begin{theoreme}[Théorème du moment cinétique]
    \[
        \left(\ddt\vect{L}_C\right)_{\text{galiléen}}
        = \vect{M}_C^{\mathrm{ext}}
    \]
\end{theoreme}

% ----------------------------------------------------------------------------
\subsection{Théorèmes de König}

\begin{theoreme}[Premier théorème de König]
    Pour tout point $A$,
    \[
        \vect{L}_A = \vect{L}_C
        + \overrightarrow{AC}\wedge m\,\vect{v}_C
    \]
\end{theoreme}

\begin{theoreme}[Deuxième théorème de König]
    L'énergie cinétique d'un corps rigide est
    \[
        T = E_K + E_K^*,
        \qquad
        E_K = \tfrac12\,m\,\|\vect{v}_C\|^2,
        \qquad
        E_K^* = \tfrac12\,\vect{\Omega}^T\mat{I}_C\,\vect{\Omega}
    \]
\end{theoreme}

\begin{preuve}
    Avec $\vect{v}_P = \vect{v}_C + \vect{\Omega}\wedge\vect{r}$,
    \begin{align*}
        T &= \frac12\int_{\mathcal{B}}\|\vect{v}_P\|^2\,\dd m
        = \frac12\int_{\mathcal{B}}
        \langle\vect{v}_C + \vect{\Omega}\wedge\vect{r},
        \vect{v}_C + \vect{\Omega}\wedge\vect{r}\rangle\,\dd m\\
        &= \frac12\,m\,\|\vect{v}_C\|^2
        + \frac12\int_{\mathcal{B}}\|\vect{\Omega}\wedge\vect{r}\|^2
        \,\dd m
        + \Bigl\langle\vect{v}_C,\ \vect{\Omega}\wedge
        \underbrace{\int_{\mathcal{B}}\vect{r}\,\dd m}_{=\,\vect{0}}
        \Bigr\rangle
    \end{align*}
    Enfin, d'après les propriétés du tenseur d'inertie,
    \[
        \frac12\int_{\mathcal{B}}\|\vect{\Omega}\wedge\vect{r}\|^2\,\dd m
        = \frac12\int_{\mathcal{B}}\|\vect{r}\wedge\vect{\Omega}\|^2\,\dd m
        = \frac12\,\vect{\Omega}^T\mat{I}_C\,\vect{\Omega}
        \qedhere
    \]
\end{preuve}

\begin{note}
    $E_K$ est l'énergie cinétique de translation : celle d'un point
    de masse $m$ placé en $C$. $E_K^*$ est l'énergie liée à la
    rotation autour du centre de masse.
\end{note}

% ----------------------------------------------------------------------------
\newpage
\subsection{Théorème de Huygens-Steiner}

Ce théorème permet de transporter le tenseur d'inertie du centre de
masse vers un autre point $A$. On pose $\vect{a} = \overrightarrow{AC}$.

\begin{theoreme}[Huygens-Steiner]
    \[
        \mat{I}_A = \mat{I}_C
        + m\bigl(\|\vect{a}\|^2\Id - \vect{a}\vect{a}^T\bigr)
    \]
\end{theoreme}

% ----------------------------------------------------------------------------
\subsection{Travail et puissance}

\begin{definition}[Puissance]
    Une force $\vect{F}$ appliquée au point matériel $P$ a pour
    \emph{puissance}
    \[
        \mathcal{P} = \vect{F}\cdot\vect{v}_P
    \]
\end{definition}

\begin{definition}[Travail]
    Le \emph{travail} de cette force entre deux instants $t_1$ et $t_2$
    est
    \[
        W = \int_{t_1}^{t_2}\mathcal{P}(t)\,\dd t
    \]
\end{definition}

Pour un solide, avec le champ des vitesses
$\vect{v}_P = \vect{v}_C + \vect{\Omega}\wedge\overrightarrow{CP}$, la
puissance totale des actions extérieures s'écrit
\[
    \mathcal{P}_{\mathrm{ext}}
    = \vect{F}_{\mathrm{ext}}\cdot\vect{v}_C
    + \vect{M}_C^{\mathrm{ext}}\cdot\vect{\Omega}
\]
\newpage
